\documentclass[
    language=english,
    doctype=article,
    institution=none,
    titlestyle=book,
]{../../omnilatex}

\RequirePackage{config/document-settings}
\addbibresource{../../bib/bibliography.bib}

% Additional packages for advanced features
\RequirePackage{amsthm}
\RequirePackage{amssymb}
\RequirePackage{algorithm2e}
\RequirePackage{tikz-cd}
\RequirePackage{makecell}
\RequirePackage{cancel}

% ─────────────────────────────────────────────────────
% Theorem Environments
% ─────────────────────────────────────────────────────
\theoremstyle{definition}
\newtheorem{theorem}{Theorem}[section]
\newtheorem{lemma}[theorem]{Lemma}
\newtheorem{proposition}[theorem]{Proposition}
\newtheorem{corollary}[theorem]{Corollary}
\theoremstyle{definition}
\newtheorem{definition}{Definition}[section]
\newtheorem{example}{Example}[section]
\theoremstyle{remark}
\newtheorem{remark}{Remark}[section]

% ─────────────────────────────────────────────────────
% Custom Commands
% ─────────────────────────────────────────────────────
\newcommand{\R}{\mathbb{R}}
\newcommand{\N}{\mathbb{N}}
\newcommand{\Z}{\mathbb{Z}}
\newcommand{\C}{\mathbb{C}}
\DeclareMathOperator*{\argmin}{arg\,min}
\DeclareMathOperator*{\argmax}{arg\,max}
\DeclareMathOperator{\Tr}{Tr}
\DeclareMathOperator{\diag}{diag}
\DeclareMathOperator{\sgn}{sgn}
\DeclareMathOperator{\Softmax}{Softmax}
\newcommand{\mat}[1]{\mathbf{#1}}
\newcommand{\tensor}[2]{#1^{#2}}
\newcommand{\hyp}{\mathcal{H}}
\newcommand{\dataset}{\mathcal{D}}
\DeclareMathOperator{\Var}{Var}
\DeclareMathOperator{\Cov}{Cov}
\newcommand{\inner}[2]{\langle #1, #2 \rangle}

\begin{document}

\maketitle

\section*{Abstract}
This article is a comprehensive showcase of advanced \LaTeX{} features available in
the OmniLaTeX template.  We demonstrate complex pgfplots visualisations, advanced
mathematical typesetting including commutative diagrams and tensor notation,
algorithm pseudocode with complexity analysis, syntax-highlighted code listings
across multiple languages, sophisticated tables with multirow and multicolumn
spanning, subfigures, cross-referencing with \verb|\autoref| and \verb|\cref|,
various citation styles, and custom TikZ diagrams including system architecture
and neural network visualisations.

\tableofcontents

% =====================================================
\section{Complex TikZ and pgfplots}
\label{sec:pgfplots}
% =====================================================

\subsection{Line and Scatter Plots with Error Bars}
\label{sec:line-scatter}

\Cref{fig:line-scatter} shows a combined line and scatter plot with error bars,
demonstrating the \texttt{error bars} library of pgfplots.

\begin{figure}[htbp]
    \centering
    \begin{tikzpicture}
    \begin{axis}[
        width=0.85\textwidth,
        height=6cm,
        xlabel={Temperature [\si{\kelvin}]},
        ylabel={Conductivity [\si{\siemens\per\metre}]},
        legend pos=north west,
        grid=major,
        thick,
        error bars/.cd,
            y dir=both,
            y explicit,
    ]
        \addplot[
            color=cBlue,
            mark=*,
            mark size=2pt,
        ] coordinates {
            (300, 12.3) +- (0, 0.8)
            (350, 18.7) +- (0, 1.2)
            (400, 24.1) +- (0, 0.9)
            (450, 31.5) +- (0, 1.5)
            (500, 38.2) +- (0, 1.1)
            (550, 42.0) +- (0, 1.8)
        };
        \addlegendentry{Measurement}

        \addplot[
            color=cRed,
            dashed,
            thick,
            domain=300:550,
            samples=50,
        ] {0.08*x - 11.5};
        \addlegendentry{Linear fit}

        \addplot[
            color=cGreen,
            mark=square*,
            mark size=2pt,
            only marks,
        ] coordinates {
            (310, 13.0)
            (370, 20.5)
            (430, 28.3)
            (480, 35.0)
            (530, 40.1)
        };
        \addlegendentry{Simulation}
    \end{axis}
    \end{tikzpicture}
    \caption{Electrical conductivity versus temperature with error bars.  Linear
    fit and simulation data are shown for comparison.}
    \label{fig:line-scatter}
\end{figure}


\subsection{Bar Chart and Grouped Plots}
\label{sec:bar-chart}

\Cref{fig:bar-chart} demonstrates grouped bar charts using the
\texttt{groupplots} library.  Two sub-panels share an axis for direct
comparison.

\begin{figure}[htbp]
    \centering
    \begin{tikzpicture}
    \begin{groupplot}[
        group style={
            group size=2 by 1,
            horizontal sep=2.5cm,
        },
        width=0.42\textwidth,
        height=6cm,
        ylabel={Accuracy [\%]},
        symbolic x coords={A, B, C, D},
        xtick=data,
        ybar=4pt,
        bar width=10pt,
        legend style={at={(0.5,-0.2)}, anchor=north, legend columns=-1,
                      font=\footnotesize},
        ymin=60, ymax=100,
        grid=major,
        grid style={dashed, gray!40},
    ]
        \nextgroupplot[title={\textbf{Dataset 1}}]
            \addplot[fill=cBlue!70] coordinates {(A,85) (B,78) (C,92) (D,88)};
            \addplot[fill=cRed!70] coordinates {(A,89) (B,83) (C,95) (D,91)};
            \addlegendentry{Baseline}
            \addlegendentry{Proposed}

        \nextgroupplot[title={\textbf{Dataset 2}}]
            \addplot[fill=cBlue!70] coordinates {(A,72) (B,68) (C,81) (D,75)};
            \addplot[fill=cRed!70] coordinates {(A,79) (B,76) (C,87) (D,82)};
            \addlegendentry{Baseline}
            \addlegendentry{Proposed}
    \end{groupplot}
    \end{tikzpicture}
    \caption{Grouped bar chart comparing baseline and proposed methods on two
    benchmark datasets across four metrics.}
    \label{fig:bar-chart}
\end{figure}


\subsection{Heatmap}
\label{sec:heatmap}

\Cref{fig:heatmap} presents a correlation heatmap generated entirely with
pgfplots, showcasing the \texttt{mesh/colormap} feature.

\begin{figure}[htbp]
    \centering
    \begin{tikzpicture}
    \begin{axis}[
        width=0.7\textwidth,
        height=7cm,
        view={0}{90},
        colormap/viridis,
        colorbar,
        colorbar style={
            ylabel={Correlation},
            ylabel style={font=\small},
        },
        xtick={0,1,2,3,4},
        xticklabels={Feature 1, Feature 2, Feature 3, Feature 4, Feature 5},
        xticklabel style={rotate=45, anchor=east, font=\small},
        ytick={0,1,2,3,4},
        yticklabels={Feature 1, Feature 2, Feature 3, Feature 4, Feature 5},
        yticklabel style={font=\small},
        title={\textbf{Feature Correlation Matrix}},
    ]
        \addplot[matrix plot*, mesh/rows=5] table {
            1.00  0.85  0.32  0.15  0.08
            0.85  1.00  0.41  0.22  0.12
            0.32  0.41  1.00  0.73  0.55
            0.15  0.22  0.73  1.00  0.68
            0.08  0.12  0.55  0.68  1.00
        };
    \end{axis}
    \end{tikzpicture}
    \caption{Heatmap of the feature correlation matrix for the five input
    variables.}
    \label{fig:heatmap}
\end{figure}


\subsection{Dual-Axis Plot}
\label{sec:dual-axis}

\Cref{fig:dual-axis} shows a dual-\(y\)-axis plot, useful when two quantities
with different units must be displayed simultaneously.

\begin{figure}[htbp]
    \centering
    \begin{tikzpicture}
    \begin{axis}[
        width=0.85\textwidth,
        height=6cm,
        xlabel={Time [\si{\second}]},
        ylabel={Voltage [\si{\volt}]},
        axis y line*=left,
        grid=major,
        thick,
        legend pos=north west,
    ]
        \addplot[cBlue, thick, domain=0:10, samples=100]
            {5*sin(deg(x*0.628)) * exp(-0.15*x)};
        \addlegendentry{Voltage}
    \end{axis}
    \begin{axis}[
        width=0.85\textwidth,
        height=6cm,
        ylabel={Current [\si{\ampere}]},
        axis y line*=right,
        axis x line=none,
        thick,
        legend pos=north east,
    ]
        \addplot[cRed, thick, dashed, domain=0:10, samples=100]
            {0.8*cos(deg(x*0.628)) * exp(-0.1*x)};
        \addlegendentry{Current}
    \end{axis}
    \end{tikzpicture}
    \caption{Dual-axis plot of voltage and current in a damped RLC circuit.}
    \label{fig:dual-axis}
\end{figure}


\subsection{Contour Plot}
\label{sec:contour}

\Cref{fig:contour} visualises a two-dimensional scalar field using contour
lines and filled regions.

\begin{figure}[htbp]
    \centering
    \begin{tikzpicture}
    \begin{axis}[
        width=0.75\textwidth,
        height=7cm,
        xlabel={\(x\)},
        ylabel={\(y\)},
        title={\textbf{Scalar Field \(f(x,y) = x^2 + y^2\)}},
        view={0}{90},
        colormap/viridis,
        colorbar,
        colorbar style={ylabel={\(f(x,y)\)}},
    ]
        \addplot3[
            contour filled={number=10},
            domain=-2:2,
            domain y=-2:2,
            samples=30,
        ] {x^2 + y^2};
        \addplot3[
            contour gnuplot={number=8},
            thick,
            domain=-2:2,
            domain y=-2:2,
            samples=30,
        ] {x^2 + y^2};
    \end{axis}
    \end{tikzpicture}
    \caption{Filled contour plot of the scalar field \(f(x,y) = x^2 + y^2\),
    demonstrating the \texttt{contour} and \texttt{colormap} features.}
    \label{fig:contour}
\end{figure}


\subsection{Pie Chart}
\label{sec:pie}

\Cref{fig:pie} shows a pie chart drawn with TikZ, illustrating the breakdown
of methods used in the benchmark study.

\begin{figure}[htbp]
    \centering
    \begin{tikzpicture}[font=\footnotesize]
        \def\radius{2.5}
        % Draw slices directly
        \draw[fill=cBlue!70,  draw=white, thick] (0,0) -- (0:\radius) arc(0:126:\radius)   -- cycle;
        \draw[fill=cRed!70,   draw=white, thick] (0,0) -- (126:\radius) arc(126:216:\radius) -- cycle;
        \draw[fill=cGreen!70, draw=white, thick] (0,0) -- (216:\radius) arc(216:288:\radius) -- cycle;
        \draw[fill=orange!70, draw=white, thick] (0,0) -- (288:\radius) arc(288:331.2:\radius) -- cycle;
        \draw[fill=violet!70, draw=white, thick] (0,0) -- (331.2:\radius) arc(331.2:360:\radius) -- cycle;
        % Percentage labels on slices
        \node[font=\bfseries] at (63:{0.65*\radius})   {35\%};
        \node[font=\bfseries] at (171:{0.65*\radius})  {25\%};
        \node[font=\bfseries] at (252:{0.65*\radius})  {20\%};
        \node[font=\bfseries] at (309.6:{0.65*\radius}) {12\%};
        \node[font=\bfseries] at (345.6:{0.65*\radius}) {8\%};
        % Legend
        \fill[cBlue!70]  (\radius+0.6,  0.0)  rectangle ++(0.3,0.3);
        \node[anchor=west] at (\radius+1.0, 0.15) {Deep Learning};
        \fill[cRed!70]   (\radius+0.6, -0.5)  rectangle ++(0.3,0.3);
        \node[anchor=west] at (\radius+1.0,-0.35) {Random Forests};
        \fill[cGreen!70] (\radius+0.6, -1.0)  rectangle ++(0.3,0.3);
        \node[anchor=west] at (\radius+1.0,-0.85) {SVMs};
        \fill[orange!70] (\radius+0.6, -1.5)  rectangle ++(0.3,0.3);
        \node[anchor=west] at (\radius+1.0,-1.35) {Grad.\ Boosting};
        \fill[violet!70] (\radius+0.6, -2.0)  rectangle ++(0.3,0.3);
        \node[anchor=west] at (\radius+1.0,-1.85) {Other};
    \end{tikzpicture}
    \caption{Market share of machine learning methods in the benchmark study,
    drawn entirely with TikZ.}
    \label{fig:pie}
\end{figure}


% =====================================================
\section{Advanced Mathematical Typesetting}
\label{sec:math}
% =====================================================

\subsection{Commutative Diagrams}
\label{sec:cd}

Commutative diagrams are typeset with the \texttt{tikz-cd} package.
\Cref{fig:cd-exact} shows an exact sequence, while \cref{fig:cd-pushout}
demonstrates a pushout square.

\begin{figure}[htbp]
    \centering
    \begin{tikzcd}[column sep=large, row sep=large]
        0 \arrow[r] &
        A \arrow[r, "f"] \arrow[d, "g"'] &
        B \arrow[r, "h"] \arrow[d, "k"'] &
        C \arrow[r] \arrow[d, "m"'] &
        0 \\
        0 \arrow[r] &
        A' \arrow[r, "f'"'] &
        B' \arrow[r, "h'"'] &
        C' \arrow[r] &
        0
    \end{tikzcd}
    \caption{A commutative diagram with exact rows.}
    \label{fig:cd-exact}
\end{figure}

\begin{figure}[htbp]
    \centering
    \begin{tikzcd}[column sep=huge, row sep=huge]
        A \arrow[r, "i_A"] \arrow[d, "i_B"'] &
        B \arrow[d, dashed, "\exists!\, u"] \\
        C \arrow[r, "j_C"'] &
        P \arrow[ul, phantom, "\ulcorner" rotate=40, very near start]
    \end{tikzcd}
    \caption{A pushout diagram in category theory.}
    \label{fig:cd-pushout}
\end{figure}


\subsection{Aligned Equations and Tensor Notation}
\label{sec:aligned}

Multi-line aligned equations are rendered with the \texttt{align} environment.
\Cref{eq:loss-decompose} decomposes the total loss.

\begin{align}
    \label{eq:loss-decompose}
    \mathcal{L}_{\text{total}}
        &= \mathcal{L}_{\text{recon}} + \beta\, \mathcal{L}_{\text{KL}}
           + \gamma\, \mathcal{L}_{\text{reg}} \\[6pt]
    \mathcal{L}_{\text{recon}}
        &= -\mathbb{E}_{q_\phi(\mat{z} \mid \mat{x})}
           \bigl[\log p_\theta(\mat{x} \mid \mat{z})\bigr] \\[4pt]
    \mathcal{L}_{\text{KL}}
        &= D_{\mathrm{KL}}\!\bigl(
            q_\phi(\mat{z} \mid \mat{x})
            \,\|\,
            p(\mat{z})
           \bigr) \\
        &= -\frac{1}{2}\sum_{j=1}^{J}
           \Bigl(
               1 + \log \sigma_j^2 - \mu_j^2 - \sigma_j^2
           \Bigr)
\end{align}

Tensor notation is used extensively in general relativity.
\Cref{eq:einstein} shows the Einstein field equations.

\begin{equation}
    \label{eq:einstein}
    R_{\mu\nu} - \frac{1}{2}\,R\,g_{\mu\nu}
        + \Lambda\, g_{\mu\nu}
        = \frac{8\pi G}{c^4}\, T_{\mu\nu}
\end{equation}

The Christoffel symbols of the second kind are defined as

\begin{equation}
    \label{eq:christoffel}
    \Gamma^{\lambda}{}_{\mu\nu}
        = \frac{1}{2}\,g^{\lambda\sigma}
          \Bigl(
              \partial_\mu g_{\nu\sigma}
              + \partial_\nu g_{\mu\sigma}
              - \partial_\sigma g_{\mu\nu}
          \Bigr)
\end{equation}


\subsection{Operator Definitions and Custom Notation}
\label{sec:operators}

We define domain-specific operators and notation.
The softmax function (\cref{eq:softmax}) and the cross-entropy loss
(\cref{eq:xent}) are fundamental in classification.

\begin{equation}
    \label{eq:softmax}
    \Softmax(z_i)
        = \frac{\exp(z_i)}{\sum_{j=1}^{K}\exp(z_j)},
        \quad i = 1, \ldots, K
\end{equation}

\begin{equation}
    \label{eq:xent}
    \mathcal{L}_{\text{CE}}
        = -\sum_{k=1}^{K} y_k \log \hat{y}_k
\end{equation}

The matrix trace satisfies \(\Tr(\mat{A}\mat{B}) = \Tr(\mat{B}\mat{A})\) and the
Hadamard product is \(\mat{A} \circ \mat{B}\).
The inner product on a Hilbert space \(\hyp\) is denoted
\(\inner{u}{v}_{\hyp}\).


\subsection{Theorem Environments}
\label{sec:theorems}

\begin{theorem}[Taylor's Theorem]
    \label{thm:taylor}
    Let \(f \colon \R \to \R\) be \(n\)-times continuously differentiable.
    For any \(x, a \in \R\),
    \begin{equation}
        f(x) = \sum_{k=0}^{n-1}
            \frac{f^{(k)}(a)}{k!}(x-a)^k
            + R_n(x),
    \end{equation}
    where the remainder satisfies
    \(R_n(x) = \frac{f^{(n)}(\xi)}{n!}(x-a)^n\) for some \(\xi\) between
    \(a\) and \(x\).
\end{theorem}

\begin{definition}[Lipschitz Continuity]
    \label{def:lipschitz}
    A function \(f \colon \R^n \to \R^m\) is \emph{Lipschitz continuous} with
    constant \(L > 0\) if
    \(\norm{f(\mat{x}) - f(\mat{y})} \le L\,\norm{\mat{x} - \mat{y}}\)
    for all \(\mat{x}, \mat{y} \in \R^n\).
\end{definition}

\begin{lemma}
    \label{lem:lipschitz-gradient}
    If \(f\) is differentiable and \(\norm{\nabla f(\mat{x})} \le L\) for all
    \(\mat{x}\), then \(f\) is Lipschitz continuous with constant \(L\).
\end{lemma}

\begin{proof}
    By the mean value theorem,
    \(\abs{f(\mat{x}) - f(\mat{y})}
        = \abs{\inner{\nabla f(\mat{\xi})}{\mat{x} - \mat{y}}}
        \le \norm{\nabla f(\mat{\xi})}\,\norm{\mat{x} - \mat{y}}
        \le L\,\norm{\mat{x} - \mat{y}}\).
\end{proof}

\begin{proposition}
    \label{prop:strongly-convex}
    If \(f\) is \(L\)-smooth and \(\mu\)-strongly convex, then
    \(\frac{\mu}{2}\norm{\mat{x} - \mat{x}^*}^2
        \le f(\mat{x}) - f(\mat{x}^*)
        \le \frac{L}{2}\norm{\mat{x} - \mat{x}^*}^2\).
\end{proposition}

\begin{corollary}
    \label{cor:condition-number}
    The condition number \(\kappa = L/\mu\) governs the convergence rate of
    gradient descent: after \(t\) iterations,
    \(f(\mat{x}_t) - f(\mat{x}^*) \le
        \frac{L\,\norm{\mat{x}_0 - \mat{x}^*}^2}{2}
        \!\left(\frac{\kappa - 1}{\kappa + 1}\right)^{\!2t}\).
\end{corollary}

\begin{example}
    \label{ex:quadratic}
    For \(f(\mat{x}) = \frac{1}{2}\mat{x}^\top\mat{Q}\mat{x}
        - \mat{b}^\top\mat{x}\) with
    \(\mu\mat{I} \preceq \mat{Q} \preceq L\mat{I}\),
    \(\kappa = \lambda_{\max}(\mat{Q})/\lambda_{\min}(\mat{Q})\).
\end{example}

\begin{remark}
    \label{rem:adaptive}
    Adaptive methods such as Adam~\cite{mckinneyPythonDataAnalysis2018} modify
    the learning rate per-parameter, effectively reducing the condition number
    in practice.
\end{remark}


% =====================================================
\section{Algorithm Pseudocode}
\label{sec:algorithm}
% =====================================================

\Cref{alg:sgd} presents stochastic gradient descent with momentum, including
complexity analysis.  The total time complexity per epoch is
\(\mathcal{O}(N\,d)\) where \(N\) is the number of samples and \(d\) is the
feature dimension.

\begin{algorithm}[htbp]
    \caption{Stochastic Gradient Descent with Momentum}
    \label{alg:sgd}
    \SetAlgoLined
    \SetKwInOut{Input}{Input}
    \SetKwInOut{Output}{Output}

    \Input{Dataset \(\dataset = \{(\mat{x}_i, y_i)\}_{i=1}^{N}\),
           learning rate \(\eta > 0\),
           momentum \(\alpha \in [0,1)\),
           epochs \(E\)}
    \Output{Optimal parameters \(\mat{\theta}^*\)}

    Initialise \(\mat{\theta}\) randomly\;
    \(\mat{v} \gets \mat{0}\) \tcp*{velocity vector}
    \For{epoch $\gets 1$ \KwTo $E$}{
        Shuffle \(\dataset\)\;
        \For{each mini-batch \(\mathcal{B} \subseteq \dataset\) of size \(B\)}{
            \(\mat{g} \gets \frac{1}{B}\displaystyle\sum_{(\mat{x}_i,y_i)\in\mathcal{B}}
                \nabla_{\mat{\theta}}\,\ell\!\bigl(f_{\mat{\theta}}(\mat{x}_i),\; y_i\bigr)\)
                \tcp*{gradient, \(\mathcal{O}(B\,d)\)}
            \(\mat{v} \gets \alpha\,\mat{v} + \eta\,\mat{g}\)
                \tcp*{velocity update, \(\mathcal{O}(d)\)}
            \(\mat{\theta} \gets \mat{\theta} - \mat{v}\)
                \tcp*{parameter update, \(\mathcal{O}(d)\)}
        }
        Evaluate loss on validation set\;
    }
    \Return \(\mat{\theta}^* \gets \mat{\theta}\)\;
\end{algorithm}


% =====================================================
\section{Code Listings}
\label{sec:listings}
% =====================================================

\subsection{Python}
\label{sec:python}

\Cref{lst:python} shows a PyTorch training loop.  The \texttt{minted} package
with the \texttt{manni} style provides syntax highlighting with line numbers.

\begin{listing}[htbp]
    \centering
    \begin{minted}[
        linenos,
        stepnumber=5,
        fontsize=\small,
        frame=leftline,
        framerule=1pt,
        rulecolor=\color{g3},
        breaklines,
    ]{python}
import torch
import torch.nn as nn
from torch.utils.data import DataLoader

def train_epoch(model, loader, optimizer, device):
    model.train()
    total_loss = 0.0
    for batch_x, batch_y in loader:
        batch_x, batch_y = batch_x.to(device), batch_y.to(device)
        optimizer.zero_grad()
        logits = model(batch_x)
        loss = nn.functional.cross_entropy(logits, batch_y)
        loss.backward()
        optimizer.step()
        total_loss += loss.item() * batch_x.size(0)
    return total_loss / len(loader.dataset)

# Hyperparameters
device = torch.device("cuda" if torch.cuda.is_available() else "cpu")
model = nn.Linear(784, 10).to(device)
optimizer = torch.optim.SGD(model.parameters(), lr=0.01, momentum=0.9)

for epoch in range(1, 11):
    loss = train_epoch(model, train_loader, optimizer, device)
    print(f"Epoch {epoch:3d}  Loss: {loss:.4f}")
    \end{minted}
    \caption{PyTorch training loop for a linear classifier on MNIST.}
    \label{lst:python}
\end{listing}


\subsection{Rust}
\label{sec:rust}

\Cref{lst:rust} demonstrates idiomatic Rust for concurrent data processing
using channels.

\begin{listing}[htbp]
    \centering
    \begin{minted}[
        linenos,
        stepnumber=5,
        fontsize=\small,
        frame=leftline,
        framerule=1pt,
        rulecolor=\color{g3},
        breaklines,
    ]{rust}
use std::sync::mpsc;
use std::thread;

fn process_data(data: Vec<f64>) -> Vec<f64> {
    let (tx, rx) = mpsc::channel();
    let chunk_size = data.len() / num_cpus::get();

    for chunk in data.chunks(chunk_size) {
        let tx = tx.clone();
        let chunk = chunk.to_vec();
        thread::spawn(move || {
            let result: Vec<f64> = chunk
                .iter()
                .map(|x| x.powi(2) + x.sqrt())
                .collect();
            tx.send(result).unwrap();
        });
    }

    drop(tx);
    rx.iter().flatten().collect()
}

fn main() {
    let data: Vec<f64> = (0..1_000_000)
        .map(|i| i as f64 * 0.001)
        .collect();
    let result = process_data(data);
    println!("Processed {} elements", result.len());
}
    \end{minted}
    \caption{Concurrent data processing in Rust using threads and channels.}
    \label{lst:rust}
\end{listing}


\subsection{SQL}
\label{sec:sql}

\Cref{lst:sql} shows a complex analytical query with window functions and
common table expressions.

\begin{listing}[htbp]
    \centering
    \begin{minted}[
        linenos,
        stepnumber=5,
        fontsize=\small,
        frame=leftline,
        framerule=1pt,
        rulecolor=\color{g3},
        breaklines,
    ]{sql}
WITH monthly_stats AS (
    SELECT
        department,
        DATE_TRUNC('month', hire_date) AS month,
        COUNT(*)                       AS new_hires,
        AVG(salary)                    AS avg_salary
    FROM employees
    WHERE hire_date >= '2024-01-01'
    GROUP BY department, DATE_TRUNC('month', hire_date)
),
ranked AS (
    SELECT
        department,
        month,
        new_hires,
        avg_salary,
        RANK() OVER (
            PARTITION BY month ORDER BY new_hires DESC
        ) AS dept_rank
    FROM monthly_stats
)
SELECT
    department,
    month,
    new_hires,
    ROUND(avg_salary, 2) AS avg_salary,
    dept_rank
FROM ranked
WHERE dept_rank <= 3
ORDER BY month DESC, dept_rank;
    \end{minted}
    \caption{SQL analytical query using CTEs and window functions.}
    \label{lst:sql}
\end{listing}


% =====================================================
\section{Complex Tables}
\label{sec:tables}
% =====================================================

\subsection{Multirow and Multicolumn Spanning}
\label{sec:multirow}

\Cref{tab:multirow} demonstrates cells that span multiple rows and columns,
a common requirement in technical reports.

\begin{table}[htbp]
    \centering
    \caption{Model performance across hardware platforms.}
    \label{tab:multirow}
    \begin{tabular}{@{}llccc@{}}
        \toprule
        \multirow{2}{*}{\textbf{Model}}
            & \multirow{2}{*}{\textbf{Backend}}
            & \multicolumn{3}{c}{\textbf{Inference Latency [\si{\milli\second}]}} \\
        \cmidrule(l){3-5}
            & & \textbf{CPU} & \textbf{GPU} & \textbf{TPU} \\
        \midrule
        \multirow{2}{*}{ResNet-50}
            & PyTorch  & 12.3 & \textbf{2.1} & 1.8 \\
            & ONNX     & 10.8 & 2.4           & --- \\
        \midrule
        \multirow{2}{*}{BERT-base}
            & PyTorch  & 45.6 & \textbf{8.3} & 6.1 \\
            & ONNX     & 38.2 & 9.7           & --- \\
        \midrule
        \multirow{2}{*}{GPT-2}
            & PyTorch  & 120.5 & \textbf{28.4} & --- \\
            & ONNX     & 98.7  & 31.2           & --- \\
        \bottomrule
    \end{tabular}
\end{table}


\subsection{Coloured Cells with \texorpdfstring{\texttt{\textbackslash cellcolor}}{cellcolor}}
\label{sec:coloured}

\Cref{tab:heatmap-table} uses background colours to highlight cells,
providing an at-a-glance performance overview.

\begin{table}[htbp]
    \centering
    \caption{Accuracy (\%) heatmap across datasets and models.}
    \label{tab:heatmap-table}
    \definecolor{heatlow}{RGB}{254,232,200}
    \definecolor{heatmid}{RGB}{253,187,45}
    \definecolor{heathigh}{RGB}{39,148,73}
    \begin{tabular}{@{}l*{5}{c}@{}}
        \toprule
        \textbf{Model}
            & \textbf{MNIST} & \textbf{CIFAR-10}
            & \textbf{ImageNet} & \textbf{COCO} & \textbf{SQuAD} \\
        \midrule
        Baseline
            & \cellcolor{heathigh}98.1
            & \cellcolor{heatmid}76.3
            & \cellcolor{heatlow}62.4
            & \cellcolor{heatlow}38.1
            & \cellcolor{heatmid}71.2 \\
        Model A
            & \cellcolor{heathigh}98.5
            & \cellcolor{heatmid}81.7
            & \cellcolor{heatmid}68.9
            & \cellcolor{heatlow}42.3
            & \cellcolor{heatmid}76.8 \\
        Model B
            & \cellcolor{heathigh}99.0
            & \cellcolor{heathigh}88.4
            & \cellcolor{heatmid}74.2
            & \cellcolor{heatmid}48.7
            & \cellcolor{heathigh}84.1 \\
        \bottomrule
    \end{tabular}
\end{table}


\subsection{Long Table Spanning Multiple Pages}
\label{sec:longtable}

\Cref{tab:long} demonstrates a table that spans multiple pages, using
\texttt{tabularray}'s \texttt{longtblr} environment with automatic
continuation headers.

\begin{longtblr}[
    caption = {Complete hyperparameter search results across all experimental
               runs.},
    label = {tab:long},
    entry = {long},
    comment = {Continued on next page},
    row{odd} = {g2!15},
]{
    colspec = {lccccr},
    row{1} = {font=\bfseries},
    hlines,
}
    \textbf{Run}
        & \textbf{Learning Rate}
        & \textbf{Batch Size}
        & \textbf{Epochs}
        & \textbf{Val. Loss}
        & \textbf{Val. Acc. (\%)}
        & \textbf{Time [\si{\minute}]} \\
    1  & 0.001  & 32   & 50  & 0.342 & 87.2 & 12  \\
    2  & 0.001  & 64   & 50  & 0.298 & 89.1 & 10  \\
    3  & 0.001  & 128  & 50  & 0.271 & 90.3 & 8   \\
    4  & 0.0005 & 32   & 80  & 0.255 & 91.0 & 18  \\
    5  & 0.0005 & 64   & 80  & 0.231 & 92.4 & 15  \\
    6  & 0.0005 & 128  & 80  & 0.218 & 93.1 & 12  \\
    7  & 0.0005 & 256  & 80  & 0.205 & 93.8 & 10  \\
    8  & 0.0002 & 64   & 120 & 0.192 & 94.2 & 22  \\
    9  & 0.0002 & 128  & 120 & 0.184 & 94.7 & 18  \\
    10 & 0.0002 & 256  & 120 & 0.178 & 95.0 & 14  \\
    11 & 0.0001 & 128  & 200 & 0.171 & 95.3 & 28  \\
    12 & 0.0001 & 256  & 200 & 0.165 & 95.6 & 22  \\
    13 & 0.0001 & 512  & 200 & 0.162 & 95.8 & 18  \\
    14 & 0.0001 & 512  & 300 & 0.158 & 96.0 & 26  \\
    15 & 0.00005& 512  & 300 & 0.155 & 96.1 & 26  \\
    16 & 0.00005& 256  & 400 & 0.153 & 96.2 & 30  \\
    17 & 0.00005& 512  & 400 & 0.151 & 96.3 & 25  \\
    18 & 0.00005& 1024 & 400 & 0.149 & 96.4 & 20  \\
    19 & 0.00005& 1024 & 500 & 0.148 & 96.5 & 24  \\
    20 & 0.00002& 1024 & 500 & 0.147 & 96.5 & 24  \\
\end{longtblr}


\subsection{Table with Footnotes and \texorpdfstring{\texttt{\textbackslash makecell}}{makecell}}
\label{sec:makecell}

\Cref{tab:footnotes} uses \texttt{makecell} for multi-line cell content and
table notes.

\begin{table}[htbp]
    \centering
    \caption{Comparison of optimiser configurations.}
    \label{tab:footnotes}
    \begin{tabular}{@{}lccl@{}}
        \toprule
        \textbf{Optimiser}
            & \textbf{\makecell[l]{Weight\\Decay}}
            & \textbf{\makecell[l]{Best\\Acc. (\%)}}
            & \textbf{Notes} \\
        \midrule
        SGD + Momentum
            & \(10^{-4}\) & 93.2
            & \makecell[l]{Standard baseline;\\requires careful LR tuning} \\
        Adam
            & ---        & 94.8
            & \makecell[l]{Fast convergence;\\may generalise poorly} \\
        AdamW
            & \(10^{-2}\) & 96.1
            & \makecell[l]{Decoupled weight decay;\\recommended for Transformers} \\
        LAMB
            & \(10^{-1}\) & 95.7
            & \makecell[l]{Designed for large batch;\scales to 64K+} \\
        \bottomrule
    \end{tabular}
    \vspace{4pt}\\
    {\footnotesize Notes: All experiments use a cosine learning rate schedule
    with warmup.  Accuracies are top-1 on ImageNet-1K.}
\end{table}


% =====================================================
\section{Figures and Subfigures}
\label{sec:figures}
% =====================================================

\subsection{Subfigures with \texorpdfstring{\texttt{\textbackslash subcaptionbox}}{subcaptionbox}}
\label{sec:subfigures}

\Cref{fig:subfigs} contains three subfigures demonstrating different plot types.

\begin{figure}[htbp]
    \centering
    \subcaptionbox{Line plot\label{fig:sf-line}}[0.3\textwidth]{
        \begin{tikzpicture}
        \begin{axis}[
            width=0.28\textwidth,
            height=4cm,
            ticks=none,
            title={\footnotesize Exponential Decay},
        ]
            \addplot[cBlue, thick, smooth, domain=0:5, samples=50]
                {exp(-0.8*x)};
        \end{axis}
        \end{tikzpicture}
    }
    \hfill
    \subcaptionbox{Bar chart\label{fig:sf-bar}}[0.3\textwidth]{
        \begin{tikzpicture}
        \begin{axis}[
            width=0.28\textwidth,
            height=4cm,
            ybar,
            bar width=12pt,
            ticks=none,
            title={\footnotesize Category Counts},
        ]
            \addplot[fill=cRed!60] coordinates {(1,8) (2,15) (3,12) (4,6)};
        \end{axis}
        \end{tikzpicture}
    }
    \hfill
    \subcaptionbox{Scatter plot\label{fig:sf-scatter}}[0.3\textwidth]{
        \begin{tikzpicture}
        \begin{axis}[
            width=0.28\textwidth,
            height=4cm,
            ticks=none,
            title={\footnotesize Clusters},
        ]
            \addplot+[only marks, mark=*, mark size=1.5pt]
                {rand*2+1};
            \addplot+[only marks, mark=square*, mark size=1.5pt]
                {rand*2+3};
        \end{axis}
        \end{tikzpicture}
    }
    \caption{Three subfigures: (\subref{fig:sf-line}) an exponential decay,
    (\subref{fig:sf-bar}) a bar chart, and (\subref{fig:sf-scatter}) a scatter
    plot.}
    \label{fig:subfigs}
\end{figure}


\subsection{TikZ Neural Network Diagram}
\label{sec:nn}

\Cref{fig:nn} shows a feed-forward neural network drawn entirely in TikZ,
illustrating the architecture visually.

\begin{figure}[htbp]
    \centering
    \begin{tikzpicture}[
        neuron/.style={
            circle,
            draw=black!80,
            thick,
            minimum size=22pt,
            inner sep=0pt,
            font=\footnotesize,
        },
        input neuron/.style={neuron, fill=cBlue!20},
        hidden neuron/.style={neuron, fill=cGreen!20},
        output neuron/.style={neuron, fill=cRed!20},
        annot/.style={text width=4em, align=center, font=\footnotesize},
    ]
        % Input layer
        \foreach \i in {1,...,4}
            \node[input neuron] (I-\i) at (0,-\i*1.2) {$x_{\i}$};

        % Hidden layer 1
        \foreach \j in {1,...,5}
            \node[hidden neuron] (H1-\j) at (2.5,-\j*1.2+0.6) {$h_{\j}$};

        % Hidden layer 2
        \foreach \j in {1,...,4}
            \node[hidden neuron] (H2-\j) at (5,-\j*1.2) {$h_{\j}'$};

        % Output layer
        \foreach \k in {1,...,2}
            \node[output neuron] (O-\k) at (7.5,-\k*1.2-0.6) {$y_{\k}$};

        % Connections: input -> hidden1
        \foreach \i in {1,...,4}
            \foreach \j in {1,...,5}
                \draw[->, thin, gray!50] (I-\i) -- (H1-\j);

        % Connections: hidden1 -> hidden2
        \foreach \j in {1,...,5}
            \foreach \k in {1,...,4}
                \draw[->, thin, gray!50] (H1-\j) -- (H2-\k);

        % Connections: hidden2 -> output
        \foreach \j in {1,...,4}
            \foreach \k in {1,...,2}
                \draw[->, thin, gray!50] (H2-\j) -- (O-\k);

        % Layer labels
        \node[annot, above=10pt] at (0,0)   {Input};
        \node[annot, above=10pt] at (2.5,0)  {Hidden 1};
        \node[annot, above=10pt] at (5,0)    {Hidden 2};
        \node[annot, above=10pt] at (7.5,0)  {Output};
    \end{tikzpicture}
    \caption{Feed-forward neural network with two hidden layers,
    4 inputs, 5 and 4 hidden units, and 2 outputs.}
    \label{fig:nn}
\end{figure}


\subsection{TikZ Circuit Diagram}
\label{sec:circuit}

\Cref{fig:circuit} uses the \texttt{circuitikz} package to draw an
RLC circuit.

\begin{figure}[htbp]
    \centering
    \begin{tikzpicture}[american]
        \draw (0,0)
            to[V, l=$V_s$, i=$i$] (0,3)
            to[R, l=$R$, i=$i_R$] (3,3)
            to[L, l=$L$, i=$i_L$] (6,3)
            to[C, l=$C$] (6,0)
            -- (0,0);
        \draw (3,3) to[short, *-] (3,2.2)
            to[short, -*] (3,1.4);
        \draw (6,3) to[short, *-] (6,2.2)
            to[short, -*] (6,1.4);
        \node at (3,1.8) {\footnotesize $v_R$};
        \node at (6,1.8) {\footnotesize $v_C$};
    \end{tikzpicture}
    \caption{Series RLC circuit showing voltage and current conventions.}
    \label{fig:circuit}
\end{figure}


\subsection{Data Visualisation with pgfplots}
\label{sec:pgfvis}

\Cref{fig:pgfvis} combines bar and line plots to show model performance
metrics.

\begin{figure}[htbp]
    \centering
    \begin{tikzpicture}
    \begin{axis}[
        width=0.85\textwidth,
        height=6cm,
        ybar,
        bar width=14pt,
        ylabel={Score},
        symbolic x coords={Precision, Recall, F1, AUC},
        xtick=data,
        ymin=0, ymax=1.05,
        legend style={at={(0.5,-0.18)}, anchor=north, legend columns=3,
                      font=\footnotesize},
        grid=major,
        grid style={dashed, gray!30},
    ]
        \addplot[fill=cBlue!60]   coordinates
            {(Precision,0.89) (Recall,0.82) (F1,0.85) (AUC,0.93)};
        \addplot[fill=cRed!60]    coordinates
            {(Precision,0.92) (Recall,0.78) (F1,0.84) (AUC,0.95)};
        \addplot[fill=cGreen!60]  coordinates
            {(Precision,0.85) (Recall,0.91) (F1,0.88) (AUC,0.91)};
        \addlegendentry{Logistic Reg.}
        \addlegendentry{Random Forest}
        \addlegendentry{Gradient Boost}
    \end{axis}
    \end{tikzpicture}
    \caption{Performance metrics for three classification models on the test
    set.}
    \label{fig:pgfvis}
\end{figure}


% =====================================================
\section{System Architecture and Flow Diagrams}
\label{sec:diagrams}
% =====================================================

\subsection{System Architecture}
\label{sec:arch}

\Cref{fig:arch} presents a high-level system architecture diagram, drawn with
TikZ using the flow chart styles from the engineering module.

\begin{figure}[htbp]
    \centering
    \begin{tikzpicture}[node distance=1.5cm, auto, >=stealth']
        % Nodes
        \node[io]       (input)     {Data Sources};
        \node[process, right of=input, xshift=1.5cm]  (ingest)  {Ingestion \& ETL};
        \node[process, right of=ingest, xshift=1.5cm] (store)   {Data Lake};
        \node[process, below of=store, yshift=-0.5cm]  (train)   {Training Pipeline};
        \node[decision, below of=train, yshift=-0.5cm] (deploy)  {Deploy?};
        \node[process, below of=deploy, yshift=-0.5cm] (serve)   {Model Serving};
        \node[io, below of=serve, yshift=-0.5cm]       (output) {Predictions};
        \node[startstop, left of=deploy, xshift=-2cm]  (log)    {Monitoring \& Logging};

        % Arrows
        \draw[arrow] (input)  -- (ingest);
        \draw[arrow] (ingest) -- (store);
        \draw[arrow] (store)  -- (train);
        \draw[arrow] (train)  -- (deploy);
        \draw[arrow] (deploy) -- node {Yes} (serve);
        \draw[arrow] (serve)  -- (output);
        \draw[arrow] (deploy) -| node [near start] {No} (log);
        \draw[arrow] (log)    |- (train);
    \end{tikzpicture}
    \caption{High-level system architecture for the ML pipeline.}
    \label{fig:arch}
\end{figure}


\subsection{Data Flow Diagram}
\label{sec:dataflow}

\Cref{fig:dataflow} shows how data moves through the processing stages.

\begin{figure}[htbp]
    \centering
    \begin{tikzpicture}[>=stealth', node distance=2cm, auto]
        \node[process, minimum width=3cm] (raw) {Raw Data};
        \node[process, right of=raw, xshift=2cm, minimum width=3cm] (clean)
            {Cleaning};
        \node[process, right of=clean, xshift=2cm, minimum width=3cm] (feat)
            {Feature Eng.};
        \node[process, right of=feat, xshift=2cm, minimum width=3cm] (model)
            {Model};

        \draw[->, thick, cBlue] (raw)  -- node {JSON, CSV} (clean);
        \draw[->, thick, cGreen] (clean) -- node {NaN fill, scale} (feat);
        \draw[->, thick, cRed]   (feat) -- node {embeddings} (model);

        % Annotation
        \draw[->, annotatearrow, thick] (2.5,-1.8) --
            node[below, annotatearrow, fill=white, inner sep=2pt, rounded corners]
            {\footnotesize Quality checks at each stage} (7,-1.8);
    \end{tikzpicture}
    \caption{Data flow from raw ingestion to model inference.}
    \label{fig:dataflow}
\end{figure}


% =====================================================
\section{Cross-References}
\label{sec:xref}
% =====================================================

OmniLaTeX provides \verb|\autoref|, \verb|\cref|, \verb|\Cref|, and
\verb|\namecref| for intelligent cross-referencing via the
\texttt{cleveref} package.  Some examples:

\begin{itemize}
    \item \verb|\autoref{eq:loss-decompile}| $\to$ \autoref{eq:loss-decompose}
    \item \verb|\cref{fig:nn}| $\to$ \cref{fig:nn}
    \item \verb|\cref{tab:multirow}| $\to$ \cref{tab:multirow}
    \item \verb|\cref{alg:sgd}| $\to$ \cref{alg:sgd}
    \item \verb|\cref{thm:taylor}| $\to$ \cref{thm:taylor}
    \item \verb|\cref{def:lipschitz}| $\to$ \cref{def:lipschitz}
    \item \verb|\cref{lst:python}| $\to$ \cref{lst:python}
    \item \verb|\cref{fig:cd-exact,fig:cd-pushout}| $\to$ \cref{fig:cd-exact,fig:cd-pushout}
    \item \verb|\namecref{sec:pgfplots}| $\to$ \namecref{sec:pgfplots}
\end{itemize}


% =====================================================
\section{Bibliography and Citations}
\label{sec:bibliography}
% =====================================================

The \texttt{biblatex} package (with \texttt{ext-authoryear} style) supports
multiple citation commands.  Examples:

\begin{itemize}
    \item \verb|\cite{einstein1905}| $\to$ \cite{einstein1905}
    \item \verb|\parencite{diracPrinciplesQuantumMechanics1981}|
          $\to$ \parencite{diracPrinciplesQuantumMechanics1981}
    \item \verb|\textcite{knuthTeXbook1986}|
          $\to$ \textcite{knuthTeXbook1986}
    \item \verb|\footcite{goossensLaTeXCompanion1993}|
          $\to$ \footcite{goossensLaTeXCompanion1993}
    \item \verb|\cites{einstein1905}{diracPrinciplesQuantumMechanics1981}|
          $\to$ \cites{einstein1905}{diracPrinciplesQuantumMechanics1981}
\end{itemize}

Multiple citations can be combined: as shown by
\textcite{einstein1905} and \textcite{knuthTeXbook1986}, foundational
works continue to shape modern research.  Further context is provided
by \parencite{baehrThermodynamikGrundlagenUnd2016,
    mollenhauerHandbuchDieselmotoren2007}.


% =====================================================
\section{Custom Commands and Operators}
\label{sec:custom}
% =====================================================

Custom commands defined in the preamble ensure consistent notation throughout.
Here is a summary of the domain-specific commands used in this article:

\begin{table}[htbp]
    \centering
    \caption{Custom commands defined for this article.}
    \label{tab:commands}
    \begin{tabular}{@{}lll@{}}
        \toprule
        \textbf{Command} & \textbf{Renders} & \textbf{Usage} \\
        \midrule
        \verb|\R|         & \(\R\)          & Real numbers \\
        \verb|\N|         & \(\N\)          & Natural numbers \\
        \verb|\mat{A}|    & \(\mat{A}\)     & Matrix / vector \\
        \verb|\inner{u}{v}| & \(\inner{u}{v}\) & Inner product \\
        \verb|\argmin|    & \(\argmin\)     & Arg minimum \\
        \verb|\argmax|    & \(\argmax\)     & Arg maximum \\
        \verb|\Tr|        & \(\Tr\)         & Matrix trace \\
        \verb|\diag|      & \(\diag\)       & Diagonal matrix \\
        \verb|\Softmax|   & \(\Softmax\)    & Softmax function \\
        \verb|\Var|       & \(\Var\)        & Variance \\
        \verb|\Cov|       & \(\Cov\)        & Covariance \\
        \bottomrule
    \end{tabular}
\end{table}

The optimisation problem from \cref{eq:loss-decompose} can be rewritten using
these operators as
\begin{equation}
    \label{eq:argmin}
    \mat{\theta}^* = \argmin_{\mat{\theta} \in \R^d}
        \frac{1}{N}\sum_{i=1}^{N}
        \ell\!\bigl(f_{\mat{\theta}}(\mat{x}_i),\; y_i\bigr)
        + \lambda\,\norm{\mat{\theta}}_2^2
\end{equation}


% =====================================================
\section{Conclusion}
\label{sec:conclusion}
% =====================================================

This article has demonstrated the breadth of advanced features available in the
OmniLaTeX template.  From complex pgfplots visualisations
(\cref{sec:pgfplots}) to commutative diagrams (\cref{sec:cd}), from algorithm
pseudocode (\cref{sec:algorithm}) to multi-language code listings
(\cref{sec:listings}), and from sophisticated tables (\cref{sec:tables}) to
TikZ system diagrams (\cref{sec:diagrams}), the template provides a
comprehensive toolkit for scientific publishing.

The combination of \texttt{cleveref} cross-references
(\cref{sec:xref}), \texttt{biblatex} citations (\cref{sec:bibliography}),
and custom mathematical commands (\cref{sec:custom}) ensures that the resulting
documents are both visually polished and semantically rich.


\printbibliography

\end{document}
